\chapter{Manufacturing: machining, casting, forming}
\label{vol08:ch06}

\epigraph{The shop floor is the truth. Drawings, BOMs, schedules, and meetings are commentary.}{anonymous manufacturing engineer}

\chapmeta{Half-life: medium-life. Archetypes: optimisation, scaling, failure. Exercise target: 40. Project track: physical.}

\noindent
Manufacturing is the discipline of converting material and design into parts. The major processes have remained recognisable for centuries (machining since the industrial revolution, casting since antiquity, forging older still); the modern engineering layer is the integration of these processes with computer-controlled equipment, quality control, supply chains, and scale. This chapter teaches the working tools: what each process does, what it costs, what it can and cannot produce, and how to choose among them. The standard reference is Kalpakjian and Schmid \cite{text:kalpakjian-manufacturing}.

\section{Manufacturing as constrained optimisation}\pathtag{core}

A manufacturer's choice of process for a part is a constrained optimisation problem: minimise cost subject to geometry (the part's shape must be producible by the process), tolerance (the process must achieve the specified accuracy), material (the process must be compatible with the part's material), volume (the process's setup cost must amortise over the production quantity), surface finish, lead time, certification.

The classical workflow is process-selection-then-design: choose the manufacturing route first, then design the part to suit. Modern \engterm{Design for Manufacturing} (DFM) and \engterm{Design for Assembly} (DFA) practices integrate these choices into the design loop: identify high-cost features at design time and remove or modify them.

\subsection{Process families}

\textbf{Subtractive (machining).} Material is removed from a workpiece by cutting, grinding, electrical discharge, or laser. Versatile, high accuracy, moderate cost per part, scales linearly in volume.

\textbf{Casting.} Molten material is poured into a mould and solidifies. Suited to complex shapes with internal features, large volumes (sand casting) or high-precision investment casting. Higher setup cost, lower per-part cost at volume.

\textbf{Forming.} Material is plastically deformed into the desired shape (forging, rolling, drawing, extrusion). Improved material properties (work hardening, grain alignment); suited to high-volume production.

\textbf{Joining (assembly).} Multiple parts combined by welding, fastening, adhesives, or other (Chapter 3).

\textbf{Additive (3D printing).} Material is added layer by layer (Chapter 7). Suited to complex internal geometry, low-volume runs, customisation.

\textbf{Surface treatment.} Material removed, added, or modified at the surface only.

The engineering choice is workload-dependent. A bracket made once for an experimental machine is machined; the same bracket made by the million is cast or formed.

\section{Machining: turning, milling, drilling, grinding}\pathtag{core}

\subsection{Turning}

A \engterm{lathe} rotates a workpiece against a stationary cutting tool. The result is a part of revolution: cylindrical, conical, threaded. The classical machine of mechanical manufacturing; modern CNC lathes (multi-axis, with live tooling) extend this to non-purely-rotational features.

The cutting parameters: spindle speed, feed rate, depth of cut. The material removal rate is $\text{MRR} = v_c f d$ where $v_c$ is cutting speed, $f$ is feed per revolution, $d$ is depth of cut. Surface finish improves with finer feed; cutting forces increase with deeper cut and harder material.

\subsection{Milling}

A \engterm{milling machine} rotates a multi-tooth cutting tool against a stationary or moving workpiece. The result is prismatic features: flats, slots, pockets, profiles. Modern CNC mills (3-axis, 5-axis) enable complex geometry; the 5-axis machines hold the part stationary while the tool reaches any face from any angle, eliminating multi-setup re-fixturing.

The cutting parameters: spindle speed, feed per tooth, depth of cut, width of cut. The chip-load (feed per tooth) drives tool life and surface finish.

\subsection{Drilling}

A \engterm{drilling machine} (or a drill on a milling machine) cuts holes with a rotating tool advanced axially. Modern hole-making includes drilling, reaming (refining drilled holes to higher accuracy), boring (enlarging a hole with a single-point tool), and tapping (cutting internal threads).

The hole diameter and tolerance are limited by the tool and the machine: drilled holes are typically $\pm 0.1$-$0.2\,\si{\milli\meter}$; reamed holes are $\pm 0.025\,\si{\milli\meter}$ or better.

\subsection{Grinding}

A \engterm{grinder} uses an abrasive wheel as the cutting tool. Used for high-accuracy finishing of hardened material (cylinders, flat surfaces, gear teeth, cutting tools themselves). Tolerances of $\pm 0.005\,\si{\milli\meter}$ and surface finishes of $R_a < 0.4\,\si{\micro\meter}$ are routine.

\subsection{Other subtractive processes}

\textbf{Electrical discharge machining (EDM).} A controlled electrical discharge erodes material. Useful for hard materials, complex geometry, internal features. Slow but accurate.

\textbf{Laser cutting.} High-power laser melts and ejects material. Used for sheet-material cutting; modern laser cutters produce parts with $R_a < 3\,\si{\micro\meter}$ on the cut edge.

\textbf{Water-jet cutting.} High-pressure water (sometimes with abrasive) cuts through material. No thermal damage; suitable for heat-sensitive materials, thick sections, composites.

\section{Casting: sand, investment, die}\pathtag{standard}

\subsection{Sand casting}

The classical process. A mould is formed in sand around a pattern (a duplicate of the part); the pattern is removed; molten metal is poured in; the metal solidifies; the sand is removed. Used for large parts, low quantities, materials that are easy to melt (iron, aluminium, brass).

Engineering virtues: low setup cost (sand patterns are cheap), broad material range, suited to very large parts (engine blocks, machine bases, marine propellers). Liabilities: rough surface finish ($R_a > 6\,\si{\micro\meter}$ typical), loose tolerances ($\pm 1$-$3\,\si{\milli\meter}$), porosity at thick sections.

\subsection{Investment casting (lost-wax)}

A wax pattern is dipped in ceramic slurry, the wax is melted out, molten metal is poured into the resulting ceramic shell. Used for high-accuracy parts, particularly aerospace turbine blades. Tolerances of $\pm 0.1\,\si{\milli\meter}$, surface finishes of $R_a < 3\,\si{\micro\meter}$ achievable. Higher per-part cost than sand casting.

\subsection{Die casting}

Molten metal is forced into a steel mould under high pressure. Used for high-volume production of aluminium, zinc, magnesium parts (automotive, consumer-electronics enclosures). Tolerances of $\pm 0.1$-$0.2\,\si{\milli\meter}$, excellent surface finish ($R_a < 1.5\,\si{\micro\meter}$). High mould cost; only economic at volumes of $10^4$ or more parts.

\subsection{Other casting}

\textbf{Permanent-mould casting.} Reusable steel mould without high-pressure injection. Intermediate between die and sand.

\textbf{Continuous casting.} Steel and aluminium produced as long slabs or bars; the working raw material of subsequent forming processes.

\textbf{Centrifugal casting.} The mould rotates; the centrifugal force fills the cavity. Used for cylindrical parts (pipes, brake drums).

\section{Forming: forging, rolling, drawing, extrusion}\pathtag{standard}

\subsection{Forging}

Plastic deformation of a heated billet between dies. Open-die forging shapes by repeated hammering (used for large parts: crankshafts, ship propellers). Closed-die (drop-forge) forging produces near-net shape in a single set of dies (connecting rods, gears, hand tools, aerospace components).

Engineering virtue: grain alignment along the load path produces parts with improved fatigue strength. Liability: limited geometry (closed-die requires draft; complex internal features impossible).

\subsection{Rolling}

Material is squeezed between rollers, reducing cross-section. Used to produce sheet, plate, structural sections (I-beams, channels, rebar). Modern continuous-rolling mills process millions of tonnes per year per facility.

\subsection{Drawing and extrusion}

\textbf{Drawing.} A wire or rod is pulled through a die, reducing its diameter. Used for fine wire, precision rod.

\textbf{Extrusion.} A billet is forced through a die, producing a constant-cross-section product. Used for aluminium and steel structural sections, plastic pipes and trim, food products.

The geometry advantage: complex cross-sections (window frames, heat-sink fins, decorative trim) can be produced economically at high volume.

\section{Sheet-metal processes}\pathtag{standard}

\subsection{Shearing and blanking}

A sheet is cut by punches and dies; the result is a blank (the desired part) and slug (waste). Used in vast volume for automotive body panels, appliances, electronics enclosures.

\subsection{Bending and forming}

A sheet is bent over a die or formed around a punch. Bending produces sheet-metal angles; forming produces deeper shapes (boxes, panels). The minimum bend radius is material-dependent (a typical rule: $r \geq t$ where $t$ is sheet thickness).

\subsection{Deep drawing}

A flat sheet is drawn into a deep cup over a punch. Used for cookware, automotive fuel tanks, beverage cans. Modern aluminium-can manufacturing produces approximately 300 cans per second per line.

\subsection{Stamping}

Combined operations (cut, bend, form) in a single press stroke or a multi-station progressive die. Used for automotive body parts.

\section{Joining: welding, brazing, soldering, mechanical}\pathtag{standard}

Chapter 3 covered the engineering of welded, adhesive, and bolted joints. This section covers the manufacturing process aspects.

\textbf{Welding} (covered in Chapter 3): high-strength permanent join with localised melting. Production welding includes robotic spot welding in automotive assembly (~$5\,000$ welds per car body), MIG and TIG for fabrications, friction stir for aluminium aerospace.

\textbf{Brazing.} A filler metal with melting point above $450\,\si{\degreeCelsius}$ but below the parent metals' is melted into the joint by capillary action. Used for plumbing fittings, automotive heat exchangers, jewellery.

\textbf{Soldering.} Same as brazing but with filler melting below $450\,\si{\degreeCelsius}$. Used for electronics PCB assembly, plumbing.

\textbf{Mechanical joining.} Riveting (aerospace, sheet-metal assemblies), screws and bolts (Chapter 3), clinching (sheet metal cold-joined without fasteners), self-pierce riveting (popular in automotive aluminium body).

\section{Surface treatments}\pathtag{standard}

A part's surface is engineered separately from its bulk. The processes:

\textbf{Heat treatment.} Annealing (soften), normalising (refine grain structure), hardening (form hard phases via quench), tempering (relieve quench stresses, balance hardness and toughness). The metallurgy was covered in Volume V.

\textbf{Case hardening.} Surface only is hardened: carburising (carbon-enriched surface), nitriding (nitrogen-enriched), induction hardening (localised heat treatment), flame hardening. Used for gears, bearings, camshafts.

\textbf{Plating.} A metal coating deposited electrolytically or by other process: chrome (decorative, corrosion-resistance), nickel (corrosion, undercoat), zinc (galvanic-corrosion protection), gold (electronics contacts).

\textbf{Anodising.} An oxide layer grown electrolytically on aluminium. Decorative and protective.

\textbf{Painting and coating.} Powder coating, e-coating, wet painting. Decorative and protective.

\textbf{Shot peening.} The surface is bombarded with small shot, introducing beneficial compressive residual stress that improves fatigue life.

\textbf{Polishing and lapping.} Mechanical refinement of surface finish.

\section{Failure: process variability, tooling wear, hidden defects}\pathtag{core}

Three classes of failure recur across manufacturing processes.

\subsection{Process variability}

A nominally well-controlled process produces a distribution of parts, not a single deterministic part. The classical engineering tools are \engterm{statistical process control} (SPC): track key measurements on each part, plot on a control chart, trigger corrective action when the measurement drifts.

\textbf{Process capability.} $C_p = (\text{USL} - \text{LSL})/(6\sigma)$, where USL and LSL are the upper and lower specification limits and $\sigma$ is the process standard deviation. A capable process has $C_p \geq 1.33$ (6$\sigma$ between specification limits, with a buffer); a Six Sigma process has $C_p \geq 2.0$.

A process with $C_p$ marginally above 1 produces a small but real defect rate; a process with $C_p \geq 2$ produces effectively no defects. The engineering response is to either tighten the process (better control, less variability) or loosen the specification (only if the function allows).

\subsection{Tooling wear}

Manufacturing tools wear in use. A drill's point dulls; a milling cutter's flutes wear; a casting mould's surface erodes; a stamping die's edge rounds. As tools wear, the parts they produce drift from nominal: holes become smaller, surfaces become rougher, edges become less sharp.

The engineering response: monitor tool wear (in-process measurements, periodic inspection); replace tools on a schedule based on the wear rate; design parts so that small tool-wear-induced drift remains within tolerance.

\subsection{Hidden defects}

Manufacturing processes can produce internal defects that visual inspection does not detect: porosity in castings, inclusions in forgings, micro-cracks in welds, residual stress from machining or heat treatment. The defects are revealed only through non-destructive testing: radiography, ultrasonics, eddy current, dye penetrant.

The engineering response is risk-based inspection: parts that are flight-critical or safety-critical get more intensive inspection than commodity parts. The aerospace and nuclear industries inspect essentially every primary-structure part; automotive inspects only critical features.

\subsection{Manufacturing capability gaps}

A specific failure mode: the part's drawing specifies features the chosen manufacturing route cannot produce. Tight tolerances on cast surfaces, sharp internal corners in machined parts, undercuts in forged parts, very thin walls in cast parts. The engineering response is design-for-manufacturing review at the drawing stage; the manufacturing engineer's input prevents costly downstream rework.

\begin{principle}[Manufacturing process determines what is possible]
A part's drawing describes geometry and tolerance; the manufacturing process determines what shapes and tolerances can be achieved economically. The engineering discipline is to make the process choice consciously, to design within its capability, and to verify (through prototype, first-article inspection, and statistical process control) that the production parts match the design.
\end{principle}

\section{Project}\pathtag{core}

\begin{project}[Manufacture a small part on actual equipment]
Machine, cast, form, or 3D-print a small part on actual equipment. Document the process and the tolerance.

\begin{enumerate}[leftmargin=*]
  \item Specify the part: a small bracket or component, $\leq 100\,\si{\milli\meter}$, with $5$-$10$ specified dimensions and $1$-$2$ critical features.
  \item Choose the manufacturing route. Justify in terms of geometry, tolerance, quantity, and material.
  \item Produce the part (manual machining, CNC machining, casting, forming, sheet-metal, or 3D printing).
  \item Measure the resulting part: at least $10$ measurements at the critical features. Compare to the drawing.
  \item Document the process: what worked, what did not, where the dimensions drifted, where the surface finish exceeded or missed the spec.
\end{enumerate}

The deliverable is the part, the measurements, and a short report (no more than five pages) on the engineering learning from the process.
\end{project}

\section{Exercises}

\subsection*{Calculation}\pathtag{core}

\begin{exercise}
For a turning operation at spindle speed $1500\,\si{rpm}$, feed $0.1\,\si{\milli\meter\per rev}$, depth $1\,\si{\milli\meter}$ on a $30\,\si{\milli\meter}$ diameter workpiece, compute the material removal rate.
\end{exercise}

\begin{exercise}
A $C_p = 1.33$ process has standard deviation $\sigma$ and specification range $\text{USL} - \text{LSL} = 0.2\,\si{\milli\meter}$. Compute $\sigma$.
\end{exercise}

\begin{exercise}
A drill produces holes with mean diameter $10.05\,\si{\milli\meter}$ and $\sigma = 0.01\,\si{\milli\meter}$. The drawing requires $10 \pm 0.05\,\si{\milli\meter}$. Compute $C_{pk}$.
\end{exercise}

\begin{exercise}
A sand casting has linear shrinkage of $1.5\%$. Compute the pattern dimension required to produce a $100\,\si{\milli\meter}$ part.
\end{exercise}

\begin{exercise}
For sheet metal of $1\,\si{\milli\meter}$ thickness, compute the bend allowance for a $90°$ bend with neutral-axis offset of $0.4$.
\end{exercise}

\begin{exercise}
Compute the press force required to shear a $2\,\si{\milli\meter}$ thick steel sheet ($\tau_y = 280\,\si{\mega\pascal}$) into a $30\,\si{\milli\meter}$ diameter blank.
\end{exercise}

\subsection*{Derivation}\pathtag{standard}

\begin{exercise}
Derive the relationship between $C_p$ and the defect rate (parts per million) for a normally distributed process.
\end{exercise}

\begin{exercise}
Show that drop-forging produces grain alignment along the deformation path. Identify the consequence for fatigue.
\end{exercise}

\begin{exercise}
Derive the minimum bend radius for sheet metal from the material's elongation limit.
\end{exercise}

\begin{exercise}
Derive the casting solidification time for a simple shape using Chvorinov's rule ($t = B(V/A)^2$).
\end{exercise}

\subsection*{Estimation}\pathtag{standard}

\begin{exercise}
Estimate the cost ratio of a single CNC-milled part vs.\ 10,000 die-cast versions of the same part.
\end{exercise}

\begin{exercise}
Estimate the tool-life of a carbide milling cutter machining mild steel at typical parameters.
\end{exercise}

\begin{exercise}
Estimate the lead time to produce a complex injection-moulded plastic part from design to first production.
\end{exercise}

\begin{exercise}
Estimate the surface finish achievable by turning vs.\ grinding the same diameter.
\end{exercise}

\begin{exercise}
Estimate the number of robotic spot welds in a typical passenger-car body.
\end{exercise}

\subsection*{Simulation}\pathtag{standard}

\begin{exercise}
Use a CAM tool to generate the G-code for a 3-axis milled pocket. Estimate machining time.
\end{exercise}

\begin{exercise}
Simulate the solidification of a casting using a finite-element tool. Identify hot spots and potential porosity.
\end{exercise}

\begin{exercise}
Implement a SPC tool: process control chart with detect-out-of-control rules.
\end{exercise}

\begin{exercise}
Simulate sheet-metal forming with deep-drawing limit-strain analysis.
\end{exercise}

\begin{exercise}
Simulate forging die-cavity-fill behaviour; identify under-fill regions.
\end{exercise}

\subsection*{Design}\pathtag{standard}

\begin{exercise}
Design the manufacturing route for a precision spindle bearing seat ($\pm 0.005\,\si{\milli\meter}$).
\end{exercise}

\begin{exercise}
Design the casting pattern for a complex engine bracket; specify draft, parting line, gating, risers.
\end{exercise}

\begin{exercise}
Design the sequence of forming operations for a sheet-metal automotive enclosure.
\end{exercise}

\begin{exercise}
Design the heat-treatment specification for a gear: case-hardened tooth, tough core.
\end{exercise}

\begin{exercise}
Design the manufacturing route for a part needing internal cooling channels (consider casting vs.\ machining vs.\ additive).
\end{exercise}

\begin{exercise}
Design the inspection plan for a die-cast part in production at $1000$ parts/day.
\end{exercise}

\begin{exercise}
Design the tooling change-out interval for a turning cutter at $1000$ parts/day with $5\%$ tool-wear-induced tolerance drift per shift.
\end{exercise}

\subsection*{Judgement}\pathtag{core}

\begin{exercise}
A team's part has machined features and cast features on the same part. State the engineering considerations and the dimensional-control challenges.
\end{exercise}

\begin{exercise}
A team uses a $C_p = 1.2$ process for a flight-critical part. State the engineering analysis.
\end{exercise}

\begin{exercise}
A team's parts pass first-article inspection but fail in production. State three engineering hypotheses.
\end{exercise}

\begin{exercise}
A team's supplier delivers parts $5\%$ out of tolerance with no apparent process change. State the engineering investigation.
\end{exercise}

\begin{exercise}
A team is choosing between machining a complex part from a billet and casting it to near-net shape then machining. State the engineering trade-offs.
\end{exercise}

\begin{exercise}
A senior engineer argues that ``modern CNC can hit any tolerance.'' Critique. State the tolerance regime where CNC requires additional process (grinding, lapping).
\end{exercise}

\begin{exercise}
A team's welded structures have inconsistent inspection results; some pass radiography, others fail. State the engineering investigation.
\end{exercise}

\begin{exercise}
A team is asked to qualify a new manufacturing supplier. State the engineering process: first-article, process qualification, sustained-monitoring.
\end{exercise}

\begin{exercise}
A team's process variability has been increasing over six months. State three engineering causes and the corresponding investigations.
\end{exercise}

\begin{exercise}
A team's part has a tight tolerance that the chosen process cannot achieve economically. State three engineering responses.
\end{exercise}

\begin{exercise}
A team's prototype part is machined; the production parts will be cast. State the engineering concerns about the transition.
\end{exercise}

\begin{exercise}
A team's surface-roughness specification has been carried over from a prior design without review. State the engineering questions to ask before accepting the specification.
\end{exercise}

\begin{exercise}
A team has the choice of one general-purpose machining centre or three specialised machines for a production line. State the engineering trade-offs and the conditions under which each choice is preferred.
\end{exercise}
